\subsection{Các ranh giới giữ context và lựa chọn kiến trúc single-outstanding}
\label{subsec:single-outstanding}

Các sơ đồ trong chương này mô tả trực tiếp mã nguồn RTL hoàn chỉnh của lõi
\texttt{MAC-F} cùng hai kiến trúc đối chứng. Cần phân biệt rõ giữa khả năng co
giãn ở cấp mô-đun con và lịch điều phối ở cấp wrapper IP: ba khối chức năng
dùng chung gồm bộ giải mã \texttt{posit\_parser\_mul} ($2$ tầng), bộ cộng
\texttt{posit\_add} ($5$ tầng) và bộ đóng gói \texttt{posit\_pack} ($2$ tầng)
được thiết kế chuẩn hóa theo giao thức bắt tay \texttt{valid/ready} co giãn đầy
đủ (\textit{elastic}) nhằm hỗ trợ ghép nối linh hoạt theo nhịp $II=1$ ở từng
tầng đơn lẻ. Tuy nhiên, ở cấp tích hợp IP \texttt{MAC-NF/F/Quire}, toàn lõi sử
dụng nhân lặp dịch--cộng tuần tự \texttt{mul\_iter\_core} ($1\text{--}12$ chu kỳ
phi đường ống) để giữ \texttt{0 DSP}, đồng thời các bài toán tích lũy chuỗi
($S_i = S_{i-1} + A_i B_i$) có phụ thuộc dữ liệu mang vòng lặp nên mỗi giao dịch
kế tiếp phụ thuộc trực tiếp kết quả $D_{i-1}$. Vì hai lý do này, wrapper chủ động
áp đặt cơ chế \textit{single-outstanding} (chỉ phát một giao dịch đang xử lý tại
một thời điểm) để tối ưu tiết kiệm logic đệm và bảo vệ ngữ cảnh $C$, chấp nhận giữ
các nhánh \textit{backpressure} liên tầng của các khối con ở trạng thái chờ tĩnh.
Độ sâu pipeline của các khối cộng và đóng gói vì vậy quyết định trực tiếp độ
trễ đơn giao dịch, không suy ra $II=1$ cho toàn IP.

\rtlfig{Ba mốc điều khiển trong \sig{mac\_f\_wrapper}. \sig{pair\_retire} chuyển payload từ Parser vào core và context C; \sig{retire} giải phóng giao dịch sau Packer.}{fig:context-control}{
\node[unit,text width=38mm] (accept) at (-5,0) {\sig{pair\_accept}\\\sig{in\_valid} $\land$ \sig{in\_ready}};
\node[unit,text width=51mm] (launch) at (0,0) {\sig{pair\_retire}\\\sig{valid\_a} $\land$ \sig{valid\_b} $\land$ \sig{valid\_c}\\$\land$ \sig{core\_ready}};
\node[unit,text width=38mm] (ret) at (5,0) {\sig{retire}\\\sig{out\_valid} $\land$ \sig{out\_ready}};
\node[ff,text width=38mm] (cfg) at (-5,-1.58) {\hw{Context cấu hình}\\\sig{n\_q[3:0]}, \sig{ops\_q}};
\node[ff,text width=51mm] (ctx) at (0,-1.58) {\hw{Context C và tích}\\frac C: 26; $sf_C$: 11; $s,Z,N$\\cờ tích, $sf_{\mathrm{core}}+1$};
\node[unit,text width=38mm] (core) at (5,-1.58) {\sig{mul\_iter\_core}\\launch, release\_result};
\draw[enable] (accept)--(cfg);
\draw[enable] (launch)--(ctx);
\draw[enable] (ret)--(core);
\draw[enable] (launch.east)--(2.83,0)|-(core.west);
\node[ff,text width=126mm] (busy) at (0,-3.15) {\sig{busy\_q} \quad
 \sig{pair\_accept}: đặt 1 \qquad|\qquad \sig{retire} hoặc reset: xóa};
\draw[enable] (accept.west)--++(-.29,0)|-(busy.west);
\draw[enable] (ret.east)--++(.29,0)|-(busy.east);
\node[tag] (gate) at (0,-4.25) {\sig{in\_ready} yêu cầu \sig{!busy\_q}, cfg hợp lệ, reset đã nhả và các khối sẵn sàng};
\draw[enable] (busy)--(gate);
}



Đầu vào cấu hình được chốt ở cạnh \texttt{pair\_accept}; các đầu vào A/B/C
vào ba Parser cùng cạnh này. \texttt{pair\_retire} là handshake của cụm
Parser với core, không phải retire ở giao diện ngoài. Ở cạnh đó, wrapper
giữ $sf_C[10:0]$, fraction $[25:0]$, dấu và metadata của C. Core nhận
X/Y đã qua OPS hoặc token bypass. Sau khi core phát \texttt{core\_done},
adapter đưa tích và context C vào S1. Wrapper chỉ nhả slot bận khi
\texttt{out\_valid} và \texttt{out\_ready} cùng bằng 1.

\begin{engineeringtable}{Chuyển quyền sở hữu dữ liệu ở các ranh giới của MAC-F.}{tab:mac-f-ownership}
\begin{tabularx}{\linewidth}{@{}>{\raggedright\arraybackslash}p{34mm}Y Y@{}}
\toprule
Sự kiện & Payload được chốt & Điều kiện giữ và giải phóng \\
\midrule
\texttt{pair\_accept} & Mã A/B/C trong Parser, n 4 bit và OPS 1 bit hợp lệ trong wrapper. & Slot bận tới retire hoặc reset. \\
\texttt{pair\_retire} & A/B đã giảm fraction, C đủ 26 bit, tổng scale và metadata đặc biệt. & C ổn định tới lúc S1 nhận tích. \\
\texttt{core\_done} & Accumulator Q24, số mũ và cờ được chuyển cùng C vào adder. & Assertion bảo vệ điều kiện nhận của adder. \\
\texttt{out\_valid} & Mã $D$ và flags 5 bit. & Giữ ổn định khi \texttt{out\_ready=0}, retire đúng một lần khi ready. \\
\bottomrule
\end{tabularx}
\end{engineeringtable}

\subsection{Backpressure, hủy giao dịch và trở lại idle}

Một giao dịch hợp lệ bắt đầu tại cạnh nhận. Nếu sink kéo ready thấp
sau đó, driver không tái cấp đầu vào đang xử lý. Trong trường hợp kết
quả đã đến Packer P2, P2 không ghi đè slot đang giữ. \texttt{out\_valid}
và payload được kiểm cùng nhau, vì giữ đúng mã nhưng hạ valid trước
handshake vẫn làm mất giao dịch. Nếu reset được assert, giao dịch
chưa retire bị hủy. Sau khi mở reset, các đồng bộ nội bộ và điều kiện
sẵn sàng chặn nhận cho đến khi pipeline có thể tiếp tục.

Trong kịch bản kiểm chứng giao thức trên bộ dữ liệu mẫu \texttt{FT-05}, phía nhận giữ \texttt{out\_ready=0} khi
kết quả xuất hiện. Ở các cạnh 14, 15 và 16 tính từ accept, \texttt{d} giữ nguyên
\texttt{1ae34800} với flags \texttt{00}. Cạnh 17 retire đúng kết quả
đó. Ca reset riêng hạ reset sau launch khi token chưa ra core,
sau đó kiểm 30 cạnh không có response cũ. Lượt sau reset dùng lại
cặp FT-05 và trả đúng mã. Trace này minh họa trực tiếp một kịch bản bắt tay
tiêu biểu ở phần thiết kế, còn thống kê diện rộng được trình bày ở Chương~\ref{ch:evaluation}.
